We highlight key findings from the development of \textbf{GenericAgent} that we believe are broadly relevant to the design of long-horizon LLM agent systems, and hope these observations will inform and inspire future research.

\textbf{Context information density is a structural constraint for all 
LLM-based agent systems.} Prior work has established that LLM reasoning 
quality does not scale linearly with context length. Information positioned 
in the middle of the context is harder to retrieve~\cite{lost_in_middle}, 
irrelevant content actively dilutes the model's attention toward 
decision-critical evidence~\cite{lost_in_middle, multi_turn_lost}, 
and the effective context window is substantially smaller than its nominal 
size~\cite{anthropic_context}. These are not artifacts of any 
particular model, but intrinsic properties of current LLM architectures. 
The implication is direct. As long as an agent uses an LLM as its reasoning 
engine, the quality of each decision step is ultimately determined within a 
single forward pass, and no amount of tooling, memory capacity, or workflow 
complexity can circumvent this constraint. Context information density is 
therefore not an optional optimization target, but a structural constraint 
that every agent system must confront by design.

\textbf{There exists a minimal complete capability set for agent systems.} 
Under the structural constraint of information density, we find that an 
agent needs to implement only three capabilities, namely \textbf{tool 
interfacing}, \textbf{context management}, and \textbf{memory formation}. 
These three capabilities correspond to the three inevitable stages in the 
agent task execution pipeline where information density is systematically 
degraded, and therefore constitute the minimal set of capabilities an agent 
framework must implement. (1) Tool interfacing is the sole channel through 
which an agent interacts with the external world. Redundant tool definitions 
consume large portions of the context budget before task execution begins, 
so interface design must be deliberately constrained to keep overhead 
minimal. (2) Context management corresponds to the input to the language 
model. Task state, intermediate results, tool outputs, and all other content 
must be actively filtered before entering the context, as loading everything 
by default allows irrelevant information to crowd out critical evidence. 
(3) Memory formation corresponds to cross-task knowledge accumulation. 
Without retaining verified content from interaction into reusable memory, 
every task begins from scratch. Any additional complexity that does not 
serve one of these three capabilities is actively degrading information 
density.


\textbf{In agent systems, lower token consumption corresponds to better 
task performance.} This finding is counterintuitive, since the prevailing 
assumption is that longer reasoning chains and more interaction turns reflect 
more thorough deliberation, and should therefore yield better outcomes. 
However, our experimental results systematically point to the opposite 
conclusion. On Lifelong AgentBench, GA consumes only 
27.7\% of Claude Code's input tokens and 15.5\% of OpenClaw's, while 
achieving a higher task completion rate of 100\%. This pattern holds 
consistently across multiple benchmarks. As discussed above, beyond a 
certain point, additional tokens do not introduce more useful information 
but instead degrade reasoning quality through positional bias, attention 
dilution, and effective window contraction~\cite{lost_in_middle, 
multi_turn_lost, anthropic_context}. An agent that consumes more tokens is more likely suffering from systematic 
failures in context management, compensating for degraded per-step decision 
quality through additional interactions rather than improving it. Token consumption reflects the symptoms of an agent's 
context management quality, rather the thoroughness of reasoning.


\textbf{Permissions define the ceiling of agent capability.} What an agent 
can perceive, what it can act upon, and what feedback it can learn from 
directly determine the complexity of reasoning chains it can develop and 
the difficulty of tasks it can solve. The scope of permissions is the 
boundary of the capability development space, which cannot be decoupled. Locking down the action 
boundary during the agent's exploration phase is equivalent to preemptively 
capping its capability ceiling at the system design stage. An agent 
restricted to reading a small set of files, unable to execute code or 
access external information, can only operate within a truncated state 
space regardless of how capable its underlying model is. Narrowing the 
exploration boundary is not a path toward building useful agents, the 
endpoint of which is a system that is safe, but useless.